%% kapa: exam=1 round=22 year=2011 period=1 form=A subject=civil_law subject_ko=민법 seq=2 number=2
%% source: Q-Net 감정평가사 기출문제 (cst00309) artlSeq=5036790 file=extracted/1008/5036790/2011년 제22회 감정평가사 1차시험 기출문제/1교시.hwp
%% conversion: kordoc 4.24.51 -> tools/gen1.py (auto); review=proofread-sonnet; flags=none
\begin{question}{2}
미성년자의 행위능력에 관한 설명으로 옳지 않은 것은?
\begin{choices}
  \item 법정대리인이 미성년자에 대하여 특정한 영업에 관하여 허락을 하면 법정대리인의 대리권은 그 범위에서 소멸한다.
  \item 법정대리인은 미성년자의 근로계약을 대리할 수 없지만, 미성년자의 임금청구는 대리할 수 있다.
  \item 상대방은 추인이 있을 때까지 미성년자의 단독행위를 거절할 수 있다.
  \item 미성년자와 계약을 체결한 자는, 계약 당시에 상대방이 미성년자임을 알았다면 그 법정대리인에게 자기의 의사표시를 철회할 수 없다.
  \item 미성년자와 계약을 체결한 자는 계약 당시에 상대방이 미성년자임을 알았더라도 그 법정대리인에게 추인 여부의 확답을 최고할 수 있다.
\end{choices}
\end{question}
\answer{2}
